% Generated draft: requires mathematical and source review.
\documentclass[11pt]{article}
\usepackage[margin=25mm]{geometry}
\usepackage{fontspec}
\setmainfont{Latin Modern Roman}
\usepackage{amsmath,amssymb,mathtools,unicode-math}
\setmathfont{Latin Modern Math}
\usepackage{enumitem,needspace,xurl,hyperref}
\setlength{\parindent}{0pt}
\setlength{\parskip}{5pt}
\setlength{\emergencystretch}{3em}
\newcommand{\N}{\mathbb{N}}
\newcommand{\Z}{\mathbb{Z}}
\newcommand{\Q}{\mathbb{Q}}
\newcommand{\R}{\mathbb{R}}
\newcommand{\C}{\mathbb{C}}
\newcommand{\F}{\mathbb{F}}
\newcommand{\A}{\mathbb{A}}
\newcommand{\PP}{\mathbb{P}}
\newcommand{\E}{\mathbb{E}}
\newcommand{\eps}{\varepsilon}
\newcommand{\dd}{\,\mathrm{d}}
\begin{document}
\section*{UnsolvedMath ID: 30001634}
\textbf{Source problem code:} catalogue entry
\subsection{Definitions and mathematical statement}
The supplied record asks about a Johnson homomorphism denoted by \(\tau_{k,\Q}\) and the group \(G_n=\mathrm{GL}(n,\Z)\). It does not specify the domain, codomain, filtration defining the homomorphism, actions of \(G_n\), or admissible values of \(n\) and \(k\).

To express the requested objects without choosing these missing data, write provisionally
\[
\tau_{k,\Q}\colon V_{n,k}\longrightarrow W_{n,k},
\]
where \(V_{n,k}\) and \(W_{n,k}\) are placeholders for the intended domain and codomain. If this is a \(G_n\)-equivariant homomorphism of rational representations, define
\[
I_{n,k}=\operatorname{im}(\tau_{k,\Q})
=\{\tau_{k,\Q}(v):v\in V_{n,k}\},
\qquad
C_{n,k}=\operatorname{coker}(\tau_{k,\Q})
=W_{n,k}/I_{n,k}.
\]
These inherit the actions \(g\cdot\tau_{k,\Q}(v)=\tau_{k,\Q}(g\cdot v)\) and \(g\cdot(w+I_{n,k})=g\cdot w+I_{n,k}\), respectively.

\textbf{Requested determination, subject to the missing specifications.} Determine the \(\mathrm{GL}(n,\Z)\)-module structure of \(I_{n,k}\); the source describes this as equivalently determining the module structure of \(C_{n,k}\). The intended range of parameters must be recovered from the cited report before this becomes a fully specified mathematical problem. The asserted equivalence also requires clarification: the exact sequence
\[
0\longrightarrow I_{n,k}\longrightarrow W_{n,k}\longrightarrow C_{n,k}\longrightarrow 0
\]
does not, by itself, make determination of the abstract module structure of the image equivalent to determination of that of the cokernel. The missing representation-theoretic hypotheses or intended meaning of determination must therefore be supplied from the source.
\subsection{Short English statement}
Determine how \(\mathrm{GL}(n,\Z)\) acts on the image and associated cokernel of the intended rational Johnson homomorphism. The supplied record leaves the homomorphism, parameter ranges, and justification of the claimed equivalence unspecified, so this is a conditional statement-only draft requiring mathematical review.
\subsection{Sources}
\begin{itemize}
\item Local UnsolvedMath registry, permanent ID \(30001634\), title \emph{Johnson-Homomorphism Images as General Linear Modules}, exact source problem code \texttt{catalogue entry}, in the Oberwolfach Reports Open Problems collection. No individual proposer is supplied; the proposed year is recorded as 2010.
\item Supplied original report citation: \emph{Teichm\"uller Theory (2011)}, Oberwolfach Reports, pp.\ 3085--3157, DOI \texttt{10.4171/owr/2010/53}. Listed also as ``original Oberwolfach problem report'': \url{https://doi.org/10.4171/owr/2010/53}. The background describes Oberwolfach Reports as documenting mathematical workshops and their open-problem discussions.
\item Supplied literature reference labeled ``recent Johnson-cokernel results'': \url{https://academic.oup.com/imrn/article/2025/10/rnaf121/8139876}. No authors or article title are supplied.
\item Literature assessment dated exactly \texttt{2026-08-21}: the supplied curation labels the problem \texttt{partially\_solved} and reports that low degrees and stable ranges of the Johnson and Andreadakis cokernels are known, while the full \(\mathrm{GL}(n,\Z)\)-module structure remains undetermined. This is supplied dated curation, not an independently verified assessment. The background expressly requests independent verification of time-sensitive status and cited claims.
\item The registry gives the statement status as \texttt{corrected\_verified} and reports that its mathematical content was checked against the official report, with notation, context, and question form normalized without changing the source-level claim. This is a supplied verification claim; no independent source consultation is asserted here.
\end{itemize}
\end{document}
